\documentclass[10pt,twocolumn]{article}

\usepackage[spanish, activeacute]{babel}
\usepackage[utf8]{inputenc}
%\usepackage{oz2tex}
\usepackage{graphicx}
\usepackage{blindtext}
\usepackage{amsmath}
\usepackage{amsfonts}
\usepackage{float}
\usepackage{url}
\usepackage{anysize} %Para usar márgenes
\usepackage{hyperref}
\marginsize{1cm}{1cm}{1.5cm}{1.5cm} %{izquierda}{derecha}{arriba}{abajo}.
\usepackage{color}
\usepackage{colortbl}
\usepackage{moreverb}

\usepackage{array}
\definecolor{naranja}{rgb}{1,0.5,0} % valores de las componentes roja, verde y azul (RGB)
\definecolor{rojo}{rgb}{1,0,0}
\definecolor{SteelBlue}{rgb}{0.3,0.5,0.7}
\definecolor{violet}{rgb}{0.35,0,0.34}

\usepackage{listings}
\usepackage{listingsutf8}
\usepackage{amsmath, amsthm, amssymb}
\lstset{ %
  language=python,
  basicstyle=\small,           % the size of the fonts that are used for the code
  numbers=none,
  numberstyle=\footnotesize,          % the size of the fonts that are used for the line-numbers
  numbersep=4pt,                  % how far the line-numbers are from the code
  backgroundcolor=\color{white},      % choose the background color. You must add \usepackage{color}
  breaklines=true,                % sets automatic line breaking
  breakatwhitespace=false,        % sets if automatic breaks should only happen at whitespace
  extendedchars=true,
  inputencoding=utf8, 
    tabsize=2,
    aboveskip=5pt,
    upquote=true,
    showstringspaces=false,
    frame=single,
    showtabs=false,
    showspaces=false,
    showstringspaces=false,
    keywordstyle=\color{blue}\ttfamily\bfseries,
    stringstyle=\color{red}\ttfamily,
    commentstyle=\color{violet}\ttfamily,
    morecomment=[l][\color{magenta}]{\#}, 
    literate={á}{{\'a}}1 {é}{{\'e}}1 {í}{{\'i}}1 {ó}{{\'o}}1 {ú}{{\'u}}1 {'}{{'}}1 {ñ}{{\~n}}1 {&}{{\&}}1 ,                                        
}

\lstset{emph={%  
    define-struct, in, then,else, letrec, symbol=,%
    },emphstyle={\color{blue}\ttfamily\bfseries}%
}%

\usepackage{tikz}
\usepackage{diagbox}

\usepackage{fancyhdr} % Para personalizar encabezados y pies de página
\usepackage{lipsum} % Para generar texto de muestra
\fancyfoot[C]{Página \thepage} % Número de página en el centro del pie de página
\renewcommand{\headrulewidth}{0pt} % Eliminar la línea del encabezado

% Aplicar el estilo de encabezado y pie de página
\pagestyle{fancy} % Para personalizar encabezados y pies de página
\usepackage{lipsum} % Para generar texto de muestra

% Definir el encabezado y pie de página
\fancyhead{} % Limpiar todos los campos de encabezado
\fancyhead[C]{
  \begin{tabular}{|c|c|c|c|}
    \hline
      \begin{tikzpicture}[baseline=(current bounding box.center)]
        \draw (0,0) circle (0.25cm);
      \end{tikzpicture} &
      \begin{tikzpicture}[baseline=(current bounding box.center)]
        \draw (0,0) circle (0.25cm);
      \end{tikzpicture} &
      \begin{tikzpicture}[baseline=(current bounding box.center)]
        \draw(0,0) circle (0.25cm);
      \end{tikzpicture} &
      \begin{tikzpicture}[baseline=(current bounding box.center)]
        \fill[black] (0,0) circle (0.25cm);
      \end{tikzpicture} \\
      \hline
  \end{tabular}} % Contenido del encabezado en la izquierda

% Aplicar el estilo de encabezado y pie de página

\title{\includegraphics[scale=0.2]{../imagenes/univalle.png} \\ Fundamentos de compilación e interpretación de lenguajes de programación \\ Primer exámen parcial, Duración 2 horas \vspace{-0.5cm}} 
\author{Carlos Andres Delgado S, Msc \\ \texttt{ carlos.andres.delgado@correounivalle.edu.co}}
\date{\vspace{-0.2cm}03 de Abril de 2025}


\newcommand{\raya}{\underline{\hspace{4cm}}}

\begin{document}

\maketitle

\thispagestyle{fancy}

\begin{enumerate}
\item {
	En un TAD las funciones que contamos son constructores, predicados y extractores que permiten:

	\begin{enumerate}
		\item Extractores: Permiten conocer que variante pertenece un tipo de dato
		\item Los predicados y extractores se denominan procedimientos observadores %correcto
		\item Constructor: Permite crear un elemento perteneciente al TAD %correcto
		\item Predicado: Permite insertar un elemento nuevo al TAD
		\item Un programador puede usar los predicados para conocer los valores internos de un TAD
	\end{enumerate}
}


\item {
  ¿Que es una función predicado?
\begin{enumerate}
  \item Es una función que construye un valor del tipo de dato.
  \item Permite extraer campos del tipo de dato
  \item Permite comparar dos datos del mismo tipo.
  \item Permite destruir un dato.
  \item Permite preguntar si un dato pertenece a un tipo de dato especifico. %correcto
\end{enumerate}
}

\item {
Indique que es cierto con respecto a un tipo abstracto de dato (TAD):
	\begin{enumerate}
		\item Si la implementación interna cambia, el programador debe ajustar sus funciones
		\item Tienen implementación e interfaz %correcto
		\item Si la implementación interna cambia, el programador no nota la diferencia %correcto
		\item Tienen Representación interna y predicados
		\item Tienen dos casos, TAD vacío que es el base y TAD extendido que permite extender a partir de una definición anterior
	\end{enumerate}

}

\item {
¿Porque se debe agregar la expresión letrec?

  \begin{enumerate}
\item Para poder declarar procedimientos, dado que en el let no podemos.
\item Dar más poder expresivo a las clausuras, al agregar más información.
\item Debido a que en el let las definiciones no pueden conocerse así mismas. %correcto
\item Porque el let no permite declarar procedimientos en su área de definiciones.
  \end{enumerate}
      
}

\item {
Un ambiente es un tipo de dato recursivo que:
	\begin{enumerate}
		\item Representa sintaxis abstracta en concreta
		\item Relaciona variables con sus valores %correcto
		\item Implementa gramáticas
		\item Representa sintaxis concreta en sintaxis abstracta
		\item Representar variables libres y ligadas
	\end{enumerate}
}

\item {
Un ambiente es un tipo de dato recursivo que:
	\begin{enumerate}
		\item El caso de ambiente vacío tiene una lista de identificadores y una lista de valores
		\item Permite representar sintaxis abstracta desde sintaxis concreta
		\item Tiene dos casos ambiente vacío y ambiente extendido %correcta
		\item Tiene tres casos: ambiente vacío, ambiente extendido y ambiente sobreextendido
		\item El caso de ambiente extendido contiene lista de identificadores, lista de valores y ambiente que extiende %correcta
	\end{enumerate}
}

\item {
Cuando cambiamos la representación interna de un tipo de dato (listas, procedimientos, bits, etc) ¿Que sucede con las funciones externas que usan la interfaz?

	\begin{enumerate}
		\item Cambian, es necesario reestructurar las funciones
		\item Cambian levemente, es necesario ajustar los condicionales de las funciones
		\item No cambian, pero el cambio de la representación puede implicar nuevas funciones auxiliares
		\item No cambian, el cambio de representación no las afecta %correcto
		\item Cambian, pero solo es necesario cambiar los condicionales y los llamados a los extractores
	\end{enumerate}
}

\item {
Dada la siguiente especificación inductiva:\\
$(0,1) \in S$  $\frac{(i,j) \in S}{(i+1,3j) \in S}$ \\ 
Marque los elementos que pertenecen a esta especificación:
\begin{enumerate}
\item $(4,81)$ %correcto
\item $(2,9)$ %correcto
\item $(2,12)$
\item $(3,0)$
\item $(1,3)$  %correcto
\end{enumerate}
}

\item {
¿Que sucede cuando cambiamos la implementación de un tipo abstracto de dato desde el punto de vista del programador (quien usa el lenguaje)?

\begin{enumerate}
\item El programador debe cambiar las funciones de la interfaz para adaptarlas al nuevo tipo de implementación
\item Las funciones de la interfaz no sirven y se hace necesario una reconstrucción del tipo de dato
\item No hay efecto, ya que el programador usa estas funciones y no trabaja directamente con la implementación %correcto
\item Se debe reestructurar todo el código, ya que la implementación exige cambiar las funciones ya realizadas por el programador.
\end{enumerate}
}

\item {
 Dada la siguiente definición de TAD
  \begin{lstlisting}
  (define-datatype dato dato?
    (datoA (x number?))
    (datoB (y symbol?) (z (list-of symbol?) (w dato?)))
  )
  \end{lstlisting}
  ¿Que declaraciones son correctas?
	\begin{enumerate}
	  \item (datoB 'x '(a b c) (datoA 5)) %Correcto
    \item (datoA 'x)
    \item (datoB 'z '(t u v) (datoB 'y '(a b c) (datoA 3))) %Correcto
    \item (datoB 'p (datoB 'x '(a b c) (datoA 3)))
    \item (datoB 'z '(a b c))
  \end{enumerate}
}

\end{enumerate}
Conteste las preguntas de 11 a 15 de acuerdo al siguiente código:

Considere el ambiente inicial env0 \texttt{(x,y,z,f), (4,2,1, closure '(a,b) +(2,+(a,b)) empty-env)} Dibuje los ambientes para la expresión:
\begin{lstlisting}
let
 a = let x = +(x,z) in (f x y)
 b = let x = +(x,y) in (f x z)
 c = let x = +(x,x) in (f y z)
 in
 letrec
  g(x,y) = if >(y,0) then +(*(2, x),(g +(x,1) -(y,1))) else let a = +(a,x) in (f a b)
  in
    ( proc (n,m) (g n m) b c)
\end{lstlisting}

\begin{enumerate}
\setcounter{enumi}{10}
\item {
El valor de a en let a = +(a,x) en el momento que es invocado es:
\begin{enumerate}
	\item 18
	\item 14
	\item 26
	\item 23 %correcto
	\item 31
\end{enumerate}
}

\item {
Para los valores de a, b, c  es cierto que:
\begin{enumerate}
	\item c = 6 %correcto
	\item a = 6
	\item c = 5 
	\item a = 9 %correcto
	\item a = 7
\end{enumerate}
}

\item {
¿Cuantos ambientes del procedimiento f se generan?
\begin{enumerate}
	\item 2
	\item 3
	\item 4 %correcto
	\item 5
	\item 6
\end{enumerate}
}

\item {
Para los valores de a, b, c  es cierto que:
\begin{enumerate}
	\item b = 6
	\item a = 3 
	\item b = 5 %correcto
	\item a = 6
	\item a = 7
\end{enumerate}
}

\item {
 En el procedimiento recursivo g(x,y) de que ambiente extiende el ambiente que contiene a en la expresión let a = +(a,x) in ...?
\begin{enumerate}
	\item Ambiente del llamado de (f a b)
	\item Ambiente inicial env0
	\item Ambiente que extiende del recursivo donde y = 1
	\item Ambiente que extiende del recursivo donde y = 0 %correcto
	\item Ambiente recursivo inicial de g(x,y)
\end{enumerate}
}

\end{enumerate}
\end{document}
